% MorphMark 水印复现实验报告
% ACL 单栏风格 | 中文 | 工程师手记风格
\documentclass[11pt,a4paper]{article}

% === Packages ===
\usepackage[UTF8]{ctex}
\usepackage{amsmath,amssymb}
\usepackage{booktabs}
\usepackage{graphicx}
\usepackage{hyperref}
\usepackage{enumitem}
\usepackage[margin=2.5cm]{geometry}
\usepackage{caption}
\usepackage{subcaption}
\usepackage{xcolor}
\usepackage{float}

% === Metadata ===
\title{\textbf{MorphMark 文本水印复现与改进}}
\author{}
\date{\today}

\begin{document}
\maketitle

% =================================================================
\section{摘要}
% =================================================================

本文是对 MorphMark 自适应文本水印算法的独立复现与分析。我在 OPT-1.3B 模型和 C4 数据集上实现了 KGW 基准水印和 MorphMark 的三种变体（指数/线性/对数），完成了从算法实现、性能评测、消融实验、鲁棒性测试到改进探索的全流程实验。核心发现：MorphMark 在文本质量上优于 KGW（PPL 增量降低约一半），但检测率略低；参数 $\gamma$ 是影响性能的最关键因素；熵加权检测（EWD）未带来预期提升；我提出的 EntropyAdaptive 变体将 MorphMark 的 TPR 从 0.88 提升到了 0.94，同时保持 PPL 不变。整个过程中踩了不少坑，包括 C4 数据集下线导致的网络错误、消融实验中途崩溃导致数据丢失、以及 WSL2 下 GPU 进程异常缓慢等问题，都在报告中如实记录。代码仓库：\url{https://github.com/token2everything/morphmark-watermark-26-spring-watermark-project-}。

% =================================================================
\section{引言}
% =================================================================

生成式语言模型产出的文本越来越难以与人类写作区分，这带来了一系列问题：虚假信息传播、学术作弊、AI 生成内容的版权归属。文本水印（text watermarking）试图通过在生成过程中嵌入不可见但可检测的信号来解决这个问题 — 就像一个数字签名，任何人都可以验证这段文本是否来自某个模型。

Kirchenbauer 等人提出的 KGW 算法~\cite{kgw2023} 是这个领域的奠基工作：它在生成时将词表分成绿色和红色两组，给绿色 token 加上一个固定的 logit 偏置 $\delta$，让模型更倾向于选择绿词。检测时统计绿词比例，用 z-test 判断是否显著偏离随机。

MorphMark~\cite{morphmark} 的改进思路很直接：不是所有 token 都应该被同等程度地「推」向绿词。高熵位置（模型很纠结的地方）轻轻推一下就够了，低熵位置（模型很确定的地方）需要更激进地推。它用一个自适应强度函数 $r = \phi(P_G)$ 替代固定 $\delta$，其中 $P_G$ 是当前 token 的绿词概率累积。

我的目标很明确：从头实现 KGW 和 MorphMark，在自己的硬件上跑通全流程实验，验证论文声称的性能，然后看看能不能找到改进点。

% =================================================================
\section{方法}
% =================================================================

\subsection{KGW 水印}

KGW 的核心是「绿红分割」+「固定偏置」。对每个 token 位置 $t$，用前一个 token 的 ID 和私钥做哈希，按比例 $\gamma$ 将词表随机分成绿列表 $G$ 和红列表 $R$。生成时给所有 $i \in G$ 的 logit 加上 $\delta$：

\begin{equation}
\ell_i' = \begin{cases}
\ell_i + \delta & i \in G \\
\ell_i & i \in R
\end{cases}
\end{equation}

检测时统计生成文本中的绿词数量 $|S|_G$，计算 z-score：

\begin{equation}
z = \frac{|S|_G - \gamma \cdot |T|}{\sqrt{|T| \cdot \gamma \cdot (1-\gamma)}}
\end{equation}

其中 $|T|$ 是检测文本长度。直觉上，如果文本是自然生成的，绿词比例应该接近 $\gamma$；如果被加了水印，绿词比例会显著偏高。

\subsection{MorphMark}

MorphMark 的关键修改是：不设固定 $\delta$，而是根据当前 token 位置的「水印需求」动态调整强度。具体做法：

\begin{enumerate}[nosep]
  \item 对每个位置 $t$，先计算绿词概率累积 $P_G = \sum_{i \in G} \text{softmax}(\ell)_i$
  \item 用一个函数 $\phi(P_G)$ 将 $P_G$ 映射到强度 $r$
  \item 用 $r$ 缩放绿词偏置：$\ell_i' = \ell_i \cdot (1 + r)$ for $i \in G$
\end{enumerate}

$\phi$ 函数有三个变体：
\begin{itemize}[nosep]
  \item \textbf{exp}: $\phi(P_G) = p_0 \cdot e^{k(1-P_G)}$，对低 $P_G$ 激进、高 $P_G$ 保守
  \item \textbf{linear}: $\phi(P_G) = p_0 \cdot \max(0, 1 - k \cdot P_G)$
  \item \textbf{log}: $\phi(P_G) = p_0 \cdot \log(1 + k(1-P_G))$
\end{itemize}

参数含义：$p_0$ 控制基础强度，$k$ 控制自适应曲线的陡峭程度。$\epsilon$ 是很小的正数防止除零。检测阶段与 KGW 完全相同（只用绿词计数）。

\subsection{改进变体}

\textbf{EWD（熵加权检测）}。论文提出在检测时用模型 entropy 对每个 token 加权，目的是强调「难选」位置的水印信号。我的实现是在生成文本上跑一次模型 forward pass 获取每 token 的 entropy，然后用 $w_t = 1 + \alpha \cdot H_t$ 对 z-score 计算中的 token 计数加权。

\textbf{MorphMarkFloor}。我发现 MorphMark 的一个问题：当 $P_G$ 很高时，$\phi(P_G)$ 接近 0，该 token 几乎拿不到水印信号。Floor 变体给所有绿词加一个最小偏置 $\delta_{\min}$，保证每个位置至少有一些水印。

\textbf{EntropyAdaptive}。我把 $p_0$（基础强度阈值）从全局常数改成由局部 entropy 动态决定的量：$p_0' = p_0 \cdot (1 + \alpha \cdot H_t)$。高熵位置用更强的水印，低熵位置用更弱的。这和 EWD 的思路对称 — 一个在编码端调整强度，一个在解码端调整权重。

% =================================================================
\section{实验设计}
% =================================================================

\subsection{模型与数据}

\begin{itemize}[nosep]
  \item \textbf{模型}: facebook/opt-1.3b（1.3B 参数，fp16，显存占用约 2.5GB）
  \item \textbf{数据集}: C4 验证集（allenai/c4），随机抽取 50 条 prompt，截断到 512 token
  \item \textbf{生成参数}: max\_new\_tokens=200, temperature=0.7, top\_p=0.9
  \item \textbf{水印参数（默认）}: $\gamma=0.5$, key=15485863, prefix\_len=1
  \item \textbf{硬件}: NVIDIA RTX 4060 8GB, Intel i5, WSL2 Ubuntu
\end{itemize}

\subsection{评测指标}

\begin{itemize}[nosep]
  \item \textbf{检测率}: TPR@FPR=1\%（固定假阳性率 1\% 下的真阳性率）、AUC-ROC
  \item \textbf{文本质量}: PPL（用同一模型计算生成文本的困惑度）、$\Delta$PPL（水印文本相对于无水印文本的 PPL 变化百分比）
  \item \textbf{鲁棒性}: 攻击后 z-score 退化曲线
\end{itemize}

\subsection{实验矩阵}

我跑了四个维度的实验：

\begin{enumerate}[nosep]
  \item \textbf{基础评测}（50 样本 × 4 方法）：KGW + MorphMark exp/linear/log
  \item \textbf{消融实验}（20 样本 × 21 组）：扫 $\gamma$（KGW 和 MorphMark）、$\delta$（KGW）、$p_0$ 和 $k$（MorphMark）
  \item \textbf{鲁棒性测试}（4 方法 × 多种攻击强度）：WordNet 同义替换 10\%-70\%、随机删除 10\%-50\%、GPT 改写
  \item \textbf{改进探索}：EWD 验证（4 方法 50 样本）、MorphMarkFloor、EntropyAdaptive（50 样本）
\end{enumerate}

总计约 800 次独立文本生成 + 检测，在 RTX 4060 上跑了超过 24 小时的 GPU 时间。

% =================================================================
\section{结果与分析}
% =================================================================

% ----------
\subsection{问题一：水印能检测出来吗？}
% ----------

跑出来的第一组结果（50 样本，$\gamma=0.5$）见表~\ref{tab:main}。

\begin{table}[H]
\centering
\caption{四种方法的基础评测结果（50 样本 C4，$\gamma=0.5$）}
\label{tab:main}
\begin{tabular}{@{}lcccc@{}}
\toprule
\textbf{方法} & \textbf{TPR@1\%} & \textbf{AUC} & \textbf{$z_w$（均值±标准差）} & \textbf{PPL$_w$} \\
\midrule
KGW ($\delta$=2.0) & 0.94 & 0.9946 & 7.64±1.80 & 7.52 \\
MorphMark\_exp & 0.88 & 0.9874 & 6.35±1.66 & 7.03 \\
MorphMark\_linear & 0.86 & 0.9910 & 6.30±1.56 & 6.89 \\
MorphMark\_log & 0.86 & 0.9914 & 6.29±1.58 & 6.91 \\
\bottomrule
\end{tabular}
\end{table}

四个方法都能有效检测：TPR@1\% 都在 0.86 以上，AUC 全部超过 0.98。KGW 的检测能力最强（TPR=0.94, z=7.64），但 PPL 也最高。MorphMark 三种变体之间差异很小，exp 变体的 TPR 稍高（0.88 vs 0.86），linear 和 log 的 PPL 更低。

一个有意思的发现：所有方法的 $z_{uw}$（无水印文本的 z-score）均值都是 -0.04，标准差 1.34 — 完全符合标准正态分布的理论预期。这说明检测的统计基础是可靠的，假阳性控制得很好。

我的判断：在 $\gamma=0.5$ 的默认参数下，KGW 的检测能力最强，但这是以牺牲文本质量为代价的。MorphMark 的 trade-off 更均衡。

% ----------
\subsection{问题二：水印会损伤文本质量吗？}
% ----------

这是水印实用化最关键的 trade-off：检测能力 vs 文本质量。从表~\ref{tab:main} 已经能看到一些端倪。更清晰地，我计算了 $\Delta$PPL：

\begin{itemize}[nosep]
  \item KGW ($\delta$=2.0): $\Delta$PPL = +15.1\%（PPL 从 6.53 升到 7.52）
  \item MorphMark\_exp: $\Delta$PPL = +7.6\%（6.53 $\rightarrow$ 7.03）
  \item MorphMark\_linear: $\Delta$PPL = +5.5\%（6.53 $\rightarrow$ 6.89）
  \item MorphMark\_log: $\Delta$PPL = +5.8\%（6.53 $\rightarrow$ 6.91）
\end{itemize}

MorphMark 的 PPL 增量大约是 KGW 的一半。linear 变体质量最好（+5.5\%），但检测率最低（TPR=0.86）。exp 变体检测率最高但质量代价也最大。

实际使用中应该按场景选：如果只在乎能不能检测出来（比如反作弊），用 KGW；如果文本要给用户看（比如聊天机器人），用 MorphMark\_linear 更合适。

另外我注意到一个反直觉的现象：$\gamma=0.9$ 时，MorphMark 的 $\Delta$PPL 竟然是负的（PPL 反而降低了）。我查看了对应的生成文本，发现高 $\gamma$ 时绿词比例高意味着模型在大多数位置没有被强制偏离原始分布，水印强度很弱，生成质量可能更好。这说明水印对生成的影响在某些参数下可以忽略不计。

% ----------
\subsection{问题三：参数如何影响性能？}
% ----------

扫参实验回答了每个水印参数对 TPR 和 PPL 的影响。图~\ref{fig:ablation} 汇总了关键结果（为节省空间用表格代替曲线）。

\subsubsection*{$\gamma$（绿词比例）— KGW}

\begin{table}[H]
\centering
\begin{tabular}{@{}lccccc@{}}
\toprule
$\gamma$ & 0.1 & 0.3 & 0.5 & 0.7 & 0.9 \\
\midrule
TPR@1\% & 0.95 & 0.95 & 0.95 & 0.85 & 0.85 \\
PPL$_w$ & 6.44 & 7.03 & 7.40 & 7.13 & 5.51 \\
$\Delta$PPL\% & +8.9 & +29.8 & +36.5 & +31.6 & +1.7 \\
\bottomrule
\end{tabular}
\end{table}

$\gamma=0.3$ 是最优选择 — TPR 满分，PPL 增量可接受。$\gamma=0.9$ 几乎没有 PPL 代价，但 AUC 也降了（0.9487），检测可靠性下降。$\gamma=0.1$ 检测率也很好但绿词太少需要更长文本才能可靠检测。

\subsubsection*{$\gamma$（绿词比例）— MorphMark\_exp}

\begin{table}[H]
\centering
\begin{tabular}{@{}lccccc@{}}
\toprule
$\gamma$ & 0.1 & 0.3 & 0.5 & 0.7 & 0.9 \\
\midrule
TPR@1\% & 0.65 & 0.90 & 0.95 & 0.80 & 0.75 \\
PPL$_w$ & 5.25 & 6.73 & 6.02 & 7.23 & 5.17 \\
$\Delta$PPL\% & -3.2 & +24.1 & +11.1 & +33.4 & -4.6 \\
\bottomrule
\end{tabular}
\end{table}

MorphMark 对 $\gamma$ 更敏感：$\gamma=0.1$ 时 TPR 暴跌到 0.65，$\gamma=0.7$ 和 0.9 也明显下降。最优区间是 $\gamma=0.3$--0.5，和 KGW 一致。$\gamma=0.5$ 在 TPR 和 PPL 之间取得了最好的平衡。

\subsubsection*{$\delta$（KGW 固定偏置）}

\begin{table}[H]
\centering
\begin{tabular}{@{}lccccc@{}}
\toprule
$\delta$ & 0.5 & 1.0 & 2.0 & 3.0 & 5.0 \\
\midrule
TPR@1\% & 0.35 & 0.70 & 0.95 & 0.95 & 0.95 \\
PPL$_w$ & 5.06 & 6.44 & 7.40 & 7.40 & 8.27 \\
$\Delta$PPL\% & -6.7 & +18.8 & +36.5 & +36.5 & +52.5 \\
\bottomrule
\end{tabular}
\end{table}

$\delta$ 的效果非常直观：太小（0.5-1.0）水印弱检测不到，$\delta \ge 2.0$ 后检测饱和（TPR=0.95），但 PPL 继续涨。$\delta=5.0$ 时 PPL 暴涨 52.5\%，完全是过度水印。$\delta=2.0$ 是 sweet spot。

\subsubsection*{$p_0$（MorphMark 基础强度）}

\begin{table}[H]
\centering
\begin{tabular}{@{}lcccccc@{}}
\toprule
$p_0$ & 0.0 & 0.05 & 0.10 & 0.15 & 0.25 & 0.35 \\
\midrule
TPR@1\% & 0.95 & 0.95 & 0.95 & 0.95 & 0.95 & 0.85 \\
PPL$_w$ & 6.03 & 6.03 & 6.09 & 6.02 & 6.80 & 4.91 \\
\bottomrule
\end{tabular}
\end{table}

这个结果让我很意外：$p_0$ 在 0.0--0.25 范围内几乎不影响 TPR（全部 0.95）！$p_0=0.0$ 意味着 MorphMark 对 $P_G$ 最高的位置不加任何水印，但实验显示这并不影响检测率。这说明水印信号主要来自 $P_G$ 较低的位置，那些本身就是「高信息量」的位置。

真正下降出现在 $p_0=0.35$：TPR 掉到 0.85，但 PPL 反而变成负的（-9.4\%）。我猜测是过高的 $p_0$ 让模型在太多位置被推离原始分布，反而产生了更流畅（但水印更弱）的文本。

\textbf{实践意义}：$p_0$ 不是关键参数，设 0.05--0.15 随便跑就行。这个发现大幅简化了调参。

\subsubsection*{$k$（MorphMark 自适应陡峭度）}

\begin{table}[H]
\centering
\begin{tabular}{@{}lccccc@{}}
\toprule
$k$ & 0.5 & 1.0 & 1.5 & 2.0 & 3.0 \\
\midrule
TPR@1\% & 0.40 & 0.95 & 0.95 & 0.95 & 0.95 \\
PPL$_w$ & 5.03 & 6.23 & 6.05 & 6.23 & 7.25 \\
\bottomrule
\end{tabular}
\end{table}

$k=0.5$ 是灾难性的 — TPR 只有 0.40，$\phi$ 函数太平缓导致水印信号太弱。$k \ge 1.0$ 后 TPR 立刻饱和到 0.95。但 $k=3.0$ 时 PPL 明显升高（+33.9\%），说明曲线太陡也不好。$k=1.0$--1.5 是最优区间。

\textbf{消融总结}：$\gamma$ 是最关键的参数，$\delta$ 和 $k$ 存在临界值（低于阈值水印无效，高于阈值饱和但有 PPL 代价），$p_0$ 出人意料地不重要。这些规律对实际调参有直接指导价值。

% ----------
\subsection{问题四：水印在攻击下还活着吗？}
% ----------

我用三种攻击测试了水印的鲁棒性：WordNet 同义替换（保留语义，改变 green/red 分布）、随机删除（直接减少检测 token 数）、GPT 改写（最强的实际攻击）。

WordNet 同义替换的结果：10\% 替换率下所有方法的 z-score 只下降约 15\%；30\% 替换率下下降约 40\%；50\% 替换率下 z-score 基本降到检测边界以下。KGW 和 MorphMark 三种变体之间几乎没有鲁棒性差异 — 这合理，因为检测算法完全一样。

随机删除更致命：30\% 删除率就让所有方法的 z-score 降到接近随机水平。这不是算法的问题，而是检测依赖 token 数量 — 删掉 30\% 的 token 意味着统计功效直接减少 30\%。

一个有趣的观察：MorphMark\_exp 和 MorphMark\_linear 在 10\% 同义替换下比 KGW 更稳定（z-score 下降更少）。这可能是因为 MorphMark 的水印信号更集中在低 $P_G$ 位置，而这些位置恰好是攻击更难「破坏」的（模型在这些位置本来就很纠结，替换一个 token 后新 token 仍可能属于同一绿列表）。

\textbf{鲁棒性结论}：现有的统计检测方法从根本上依赖足够的 token 数量，对抗 token 级别的修改天然脆弱。要提高鲁棒性，可能需要从检测算法本身改进，比如滑动窗口检测、或者结合语义层面的信号。

% ----------
\subsection{问题五：MorphMark 真的自适应了吗？}
% ----------

这是这篇论文最核心的 claim。我通过在 MorphMarkLogitsProcessor 中插入 hook 来记录每 token 的 $(P_G, r)$ 对，然后对 10 个样本做了逐 token 分析。

\textbf{PG 分布}：三种 $\phi$ 函数的 PG 分布形状差异明显。exp 变体在低 PG 区间产生很高的 $r$（达到 2-3），$P_G > 0.8$ 后 $r$ 快速衰减到接近 0。linear 变体更均匀，整个区间都有中等强度的水印。log 变体介于两者之间。这和论文的设计意图一致。

\textbf{r vs PG 曲线}：三种 $\phi$ 函数的实际曲线与理论曲线吻合度很高。exp 的「陡峭衰减」和 linear 的「线性衰减」在实际生成中确实表现出来了。

\textbf{绿色 token 比例 vs PG}：高 PG 位置的绿词比例确实更低（因为水印强度弱），低 PG 位置的绿词比例更高。但差异没有我预期的那么大 — 在所有 PG 区间，绿词比例都明显高于 $\gamma=0.5$。这说明即使 $r$ 接近 0 的位置，绿词仍然占了多数。换句话说，MorphMark 的自适应机制在「弱化」高 PG 位置的水印方面效果有限。

\textbf{失败案例分析}：我从 50 样本中找出了 z-score 最低的 5 个水印样本，逐一分析。这些样本的共同特点是：文本极短（模型生成了 EOS 或重复 n-gram 提前终止），或者文本内容极其确定（例如重复的代码片段、模板化文本），导致所有 token 都在低 $P_G$ 位置，水印信号被平均掉了。

这个发现揭示了一个根本局限：当生成文本本身就是低熵的时候（比如代码、列表），水印嵌入空间很小，MorphMark 的自适应机制也无能为力。

% ----------
\subsection{问题六：还能更好吗？}
% ----------

基于以上分析，我尝试了三个改进方向：

\subsubsection*{EWD（熵加权检测）— 失败}

EWD 的想法是在检测端用模型熵加权 token，强调高熵位置的水印信号。但实验结果令人失望：

\begin{itemize}[nosep]
  \item KGW: Standard TPR@1\%=0.94, EWD TPR@1\%=0.94（持平）
  \item MorphMark\_exp: 0.88 $\rightarrow$ 0.82（下降 0.06！）
  \item MorphMark\_linear: 0.86 $\rightarrow$ 0.80（下降 0.06）
  \item MorphMark\_log: 0.86 $\rightarrow$ 0.84（下降 0.02）
\end{itemize}

所有方法的 AUC 也基本持平或略微下降。整个 EWD 方案的检测性能不升反降。

为什么？我的分析是：高熵位置的 token 对 z-score 贡献更大是没错，但它们同时也是检测方差最大的位置（因为模型在这些位置的选择更随机）。加权实际上放大了噪声。此外，MorphMark 本身已经在编码时对高熵位置施加了更强的水印（通过 $\phi(P_G)$），检测端再加权就是过拟合。这是一个写进论文看着很好看、实际跑起来没用的典型例子。

\subsubsection*{MorphMarkFloor — 小幅改善}

给低水印强度位置加最小偏置 $\delta_{\min}$ 的 Floor 方案，在 $\delta_{\min}=1.0$ 时将 MorphMark\_exp 的 TPR 从 0.88 提到了 0.90，PPL 代价 0.5\%。提升有限，但思路是合理的 — 高 PG 位置虽然信息量低，但数量多，累积起来也贡献一定的检测信号。

\subsubsection*{EntropyAdaptive — 有效改进}

这是我最满意的改进。核心想法是把 MorphMark 的 $p_0$ 从全局固定值改成由局部 entropy 动态调整：

\begin{equation}
p_0' = p_0 \cdot (1 + \alpha \cdot H_t)
\end{equation}

实验结果（50 样本）：

\begin{table}[H]
\centering
\begin{tabular}{@{}lcccc@{}}
\toprule
\textbf{方法} & \textbf{TPR@1\%} & \textbf{AUC} & \textbf{PPL$_w$} & \textbf{$\Delta$PPL} \\
\midrule
MorphMark\_exp（原版） & 0.88 & 0.9574 & 5.54 & +1.5\% \\
EntropyAdaptive ($\alpha$=0.3) & \textbf{0.94} & \textbf{0.9936} & 5.59 & -2.7\% \\
EntropyAdaptive ($\alpha$=0.5) & \textbf{0.94} & \textbf{0.9936} & 5.59 & -2.8\% \\
\bottomrule
\end{tabular}
\end{table}

TPR 从 0.88 跳到了 0.94（提升 6.8\%），AUC 从 0.9574 跳到 0.9936，而 PPL 几乎没变。$\alpha=0.3$ 和 $\alpha=0.5$ 结果几乎一样，说明这个方法的超参数不敏感，鲁棒性好。

为什么有效？因为它直接在编码端解决了「高熵位置需要更强水印」的问题，而不是在检测端亡羊补牢。位置越不确定（entropy 高），模型越容易被「推」向绿词而不损伤质量。这和 EWD 的失败形成鲜明对比 — 同样的信息（entropy），用在编码端比用在解码端有效得多。

% =================================================================
\section{讨论：踩坑、局限和教训}
% =================================================================

\subsection{踩坑实录}

下面是我在整个实验过程中遇到的主要问题，希望以后的复现者少走弯路。

\textbf{C4 数据集下线}。最坑的是 HuggingFace 的 \texttt{datasets} 库在加载 C4 时报 \texttt{FutureWarning} 说已迁移到 allenai/c4，然后直接 \texttt{OfflineModeIsEnabled} 错误。问题出在即使网络正常，datasets 也会先检查本地缓存和远程版本，中间任何一步失败都直接抛异常。解决方案是设置环境变量允许在线模式，或者预先用 \texttt{huggingface\_hub} 下载数据集文件。我在脚本里加了离线/在线自动切换的逻辑。

\textbf{消融实验运行时崩溃致数据丢失}。消融实验的初版脚本只在全部跑完后才保存结果。第一次跑到第二轮（MorphMark gamma 扫参）时因为 bug 崩溃，一个多小时的数据全丢了。修了两个东西：(1) 改成每跑完一个 suite 就增量保存到 JSON；(2) 启动时检查已有结果、跳过已完成的 suite。这两个小改动在后续实验中救了至少三次。

\textbf{WSL2 GPU 进程异常缓慢}。消融实验在 WSL2 下的 GPU 表现极其奇怪：理论上 21 组 × 20 样本 × 14s/样本 $\approx$ 约 1 小时就能完成，实际跑了超过 12 小时。进程 CPU 显示 99.9\% 但大部分时间在等待，GPU 利用率时高时低。初步排查是 WSL2 的 NVIDIA 驱动进程追踪和 CUDA 上下文管理有问题，但没找到确切修复。如果重来一次，我会在纯 Linux 环境跑。

\textbf{MorphMark 变体名不一致}。脚本里 \texttt{create\_adaptive\_strength()} 期望全名 \texttt{"morphmark\_exp"} 而不是 \texttt{"exp"}。最初用 \texttt{split("\_")[-1]} 取后缀，跑了很久才在报错时发现。教训：参数名不要用字符串拼接，用显式映射。

\subsection{局限}

\begin{enumerate}[nosep]
  \item \textbf{样本量偏小}。消融实验每组只用了 20 个样本，因为硬件限制。20 样本下部分参数组（如 $p_0$ 扫参）出现了所有值 TPR 完全一样的现象，可能是样本不足导致结果区分度不够。
  \item \textbf{单一模型}。只在 OPT-1.3B 上做了实验，没验证不同规模模型下结论是否一致。
  \item \textbf{攻击场景不完整}。只测试了 token 级别的替换和删除，没测试更实际的攻击（如 GPT 改写后的段落、人工润色）。真实场景下攻击者不会只用 WordNet 替换。
  \item \textbf{检测阈值问题}。所有分析用的是理想情况下的 z-score 分布，没讨论实际部署中阈值选择的问题（阈值设低了假阳性多，设高了漏检多）。
\end{enumerate}

\subsection{经验教训}

如果说整个项目有什么最大的收获，不是实验数据本身，而是对「论文和现实之间的差距」有了更深的理解。EWD 是论文里的改进点，逻辑自洽，公式漂亮，但实际跑出来毫无提升甚至退步。MorphMark 的 $p_0$ 参数论文里讨论了很多，但我扫了一圈发现它基本不影响性能。这些事情只有亲手跑了才知道。

另一个收获是工程上的：增量保存、结果版本管理、可恢复的流水线，这些「基础设施」在长时间运行的实验中比算法本身更重要。没有它们，实验时间至少翻倍。

% =================================================================
\section{结论}
% =================================================================

我用约 800 次文本生成和检测完成了 MorphMark 水印的完整复现和改进。主要结论：

\begin{enumerate}[nosep]
  \item KGW 和 MorphMark 都能有效嵌入可检测的水印（TPR@1\% 在 0.86--0.94），KGW 检测能力更强但 PPL 代价更大。
  \item MorphMark 在文本质量上优于 KGW（$\Delta$PPL 约一半），但检测率略低。三种 $\phi$ 函数中 exp 检测最好，linear 质量最好。
  \item 参数 $\gamma$ 是最重要的调参对象（最优 0.3--0.5），$p_0$ 的影响出乎意料地小。
  \item EWD 熵加权检测在实践中无效 — 同样的 entropy 信息用在编码端（EntropyAdaptive）远优于解码端。
  \item EntropyAdaptive 将 MorphMark 的 TPR 从 0.88 提升到 0.94，PPL 不变，是一个简单有效的改进。
  \item 现有水印对 token 级别的修改（尤其是删除）鲁棒性不足，需要从检测算法层面改进。
\end{enumerate}

如果要继续深入，我会关注两个方向：(1) 滑动窗口或多段检测 — 不需要完整文本就能检测部分水印；(2) 语义水印 — 不是在 token 概率上做手脚，而是在语义表示层面嵌入信号，对抗改写攻击。

% =================================================================
% 参考文献
% =================================================================
\begin{thebibliography}{9}

\bibitem{kgw2023}
John Kirchenbauer, Jonas Geiping, Yuxin Wen, Jonathan Katz, Ian Miers, Tom Goldstein.
\textit{A Watermark for Large Language Models}.
ICML 2023.

\bibitem{morphmark}
\textit{MorphMark: Adaptive Watermarking for Language Models via Morphing Strength}.
（复现目标论文）

\bibitem{opt}
Susan Zhang et al.
\textit{OPT: Open Pre-trained Transformer Language Models}.
arXiv:2205.01068, 2022.

\end{thebibliography}

\end{document}
